\section[Adoption]{A realistic starting point}
\sectionframe{08}{A realistic starting point}
  {What should I change in my repository first?}
  {\faSeedling}

% ---------------------------------------------------- other skill packages
% \pkgcard{k}{owner}{repository}{head colour}{author}{what it is}{its steps}{its rule for tests}: card k
% (0 to 2), 46 mm wide and 3 mm apart; the green box at its foot holds the rule, level with the other
% cards' rules
\newcommand\pkgcard[8]{%
  \pgfmathsetmacro\pkx{#1*49}%
  \begin{scope}
    \clip[rounded corners=1.2mm] (\pkx,0) rectangle ++(46,-45.4);
    \fill[white] (\pkx,0) rectangle ++(46,-45.4);
    \fill[#4] (\pkx,0) rectangle ++(46,-9.2);
    \fill[greentint] (\pkx,-32.6) rectangle ++(46,-12.8);
    \fill[testgreen] (\pkx,-32.6) rectangle ++(1.2,-12.8);
  \end{scope}
  \draw[udsgreylight,line width=0.5pt,rounded corners=1.2mm] (\pkx,0) rectangle ++(46,-45.4);
  \node[anchor=north west,inner sep=0pt,font=\ttfamily\fontsize{7}{8.4}\selectfont,text=white!70!#4,
    align=left] at (\pkx+1.8,-1.5) {#2/\\\textcolor{white}{\bfseries #3}};
  \node[pkg who] at (\pkx+1.8,-10.8) {#5};
  \node[pkg what] at (\pkx+1.8,-14.8) {#6};
  \node[pkg text] at (\pkx+1.8,-22.6) {#7};
  \node[pkg text,text width=40.6mm] at (\pkx+3.6,-34.4) {\textcolor{testgreen}{\faCheck}\enspace #8};}
\begin{frame}[source={GitHub, 3 Nov 2026: obra/superpowers v7.0.1 · multica-ai/andrej-karpathy-skills · mattpocock/skills v1.3.1}]{Other skill packages}
  \begin{tikzpicture}[x=1mm,y=1mm,
      pkg who/.style={anchor=north west,inner sep=0pt,font=\scriptsize,text=udsgrey,text width=42.4mm,align=left},
      pkg what/.style={anchor=north west,inner sep=0pt,font=\footnotesize\bfseries,text=udsdark,
        text width=42.4mm,align=left,execute at begin node={\hyphenpenalty=10000\raggedright}},
      pkg text/.style={anchor=north west,inner sep=0pt,font=\scriptsize,text=udsdark,text width=42.4mm,
        align=left,execute at begin node={\hyphenpenalty=10000\raggedright}}]
    \useasboundingbox (0,0) rectangle (144,-54);
    \pkgcard{0}{mattpocock}{skills}{udsdark}{Matt Pocock}{small skills you combine}
      {grill, spec, tickets, implement, review}{\fn{/tdd}: red before green, at agreed seams}
    \pkgcard{1}{obra}{superpowers}{udsblue}{Jesse Vincent}{a whole method in 15 skills}
      {brainstorm, plan, execute, review, finish}{no production code without a failing test first}
    \pkgcard{2}{multica-ai}{andrej-karpathy-skills}{udsblue}{from Karpathy's notes}{one \fn{CLAUDE.md}, four rules}
      {think first, keep it simple, change only what is asked, set checkable goals}
      {a bug fix starts with a test that reproduces the bug}
    % what all three share
    \fill[greentint] (0,-48) rectangle (144,-54);
    \node[anchor=west,inner sep=0pt,font=\footnotesize\bfseries,text=testgreen] at (2,-51)
      {\faCheckDouble\enspace in all three, the check is set before the code is written};
  \end{tikzpicture}
  \note{\textbf{37:00--38:00}\par
    Pocock's collection is one of several; two others are worth a look.\par
    Superpowers, by Jesse Vincent, is a whole method, from brainstorming to finishing the
      branch; its test-driven skill allows no production code without a failing test.\par
    The Karpathy skills are one CLAUDE.md with four rules, written by Jiayuan Zhang from Andrej
      Karpathy's notes on how models go wrong in code; the last rule turns a task into a goal
      that can be checked, such as a test that reproduces a bug.\par
    They differ in how much of the process they take over; Pocock's README prefers small skills
      you combine to systems that own the process.\par
    What they share is section 6: the check is decided before the code.\par
    As with any skill, read them before installing.}
\end{frame}

% --------------------------------------------------------------- checklist
\begin{frame}{Checklist for your repository}
  \setbeamertemplate{itemize item}{\raisebox{-0.1ex}{\textcolor{udsblue}{\faSquare[regular]}}}
  \setlength\leftmargini{5mm}
  \begin{columns}[T,onlytextwidth]
    \begin{column}{0.3\textwidth}
      \textbf{\color{udsdark}Groundwork}\par\vskip1mm
      \begin{itemize}\footnotesize\setlength\itemsep{2mm}
        \item a git repository; commit before each agent task
        \item the Python environment (conda, uv, etc.)
        \item Ruff and mypy configured, run after every edit
      \end{itemize}
    \end{column}
    \begin{column}{0.31\textwidth}
      \textbf{\color{udsblue}Context}\par\vskip1mm
      \begin{itemize}\footnotesize\setlength\itemsep{2mm}
        \item \fn{AGENTS.md}: commands, boundaries, pointers
        \item \fn{GLOSSARY.md}: the key terms
        \item \fn{CODING\_STANDARDS.md}: what review checks
        \item \fn{docs/adr/}: decisions people reopen
      \end{itemize}
    \end{column}
    \begin{column}{0.33\textwidth}
      \textbf{\color{udsred}Workflow}\par\vskip1mm
      \begin{itemize}\footnotesize\setlength\itemsep{2mm}
        \item a workflow that suits you, e.g.\
          \href{https://github.com/mattpocock/skills}{\faGithub\enspace\underline{\fn{mattpocock/skills}}}
        \item your own skills for the tasks you repeat
        \item tests agreed before the code
        \item one worktree, or a forked session, per task
        \item you review the code and the science
      \end{itemize}
    \end{column}
  \end{columns}
  \note{\textbf{38:00--38:50}\par
    Groundwork: a git repository with a commit before every agent task, the Python environment
      with conda or uv, and Ruff and mypy run after every edit.\par
    Context: AGENTS.md, the glossary, the coding standards and ADRs, so that requests can stay
      short.\par
    Workflow: choose one that suits you, Pocock's skills for example; write your own skills for
      the tasks you repeat; agree the tests before the code; give each task its own worktree or
      forked session; and review the code and the science yourself.\par
    None of this needs a new tool, and most of it helps human collaborators too.}
\end{frame}

% --------------------------------------------------------------- thank you
% the title slide's layout: its blue panel on the right, the thanks where the title was; drawn
% inside the frame, so that the notes' thumbnail shows it too
\begin{frame}[plain,noframenumbering]
\vbox to 0pt{\vskip-2mm\hbox to 0pt{\hskip-8mm\begin{tikzpicture}[x=1mm,y=1mm]
  \useasboundingbox (0,0) rectangle (160,-90);
  \titlepanel
  \node[anchor=north west,inner sep=0pt] at (8,-8) {\includegraphics[width=30mm]{\udslogo}};
  \node[anchor=north west,inner sep=0pt,text width=88mm,align=left] at (8,-36) {%
    {\fontsize{32}{36}\selectfont\bfseries\color{udsblue}Thank you\par}\vskip4mm
    \textcolor{udsred}{\rule{14mm}{0.8pt}}\par\vskip3mm
    {\fontsize{11}{13}\selectfont\color{udsdark}\insertauthor\par}};
\end{tikzpicture}\hss}\vss}
\par\rule{0pt}{80mm}
\note{\textbf{38:50--39:20}\par
    Thank you.\par
    Questions to open the discussion: which of your projects is not under git yet; which
      repository you would try this on first, and who reviews it; which decision in your project deserves an
      ADR; which single command tells you that your code still works.\par
    Templates, tool details and references follow in the appendix.}
\end{frame}
